\chapter{Introduction}

Streaming catalogs currently host over 250 million recordings, with 100,000 new tracks uploaded every single day~\cite{ref:ifpi2026}. While casual listeners enjoy endless background music, independent music producers hunting for samples face an exhausting search process. Audio sampling underpins roughly 17\% of Billboard Hot 100 hits~\cite{ref:tracklib2022}. Yet commercial streaming platforms treat every song as an undifferentiated, continuous audio file. Producers rarely sample entire songs. Instead, they search for localized 4-bar drum breaks, clean keyboard progressions, or solo intros without vocals. 

Our project, \textbf{"Better Spotify / Sample Browser"} (hereinafter referred to as \textbf{Sample Browser}), tackles this bottleneck directly. More than half of active music makers work regular jobs outside the music industry~\cite{ref:ukcensus2023}. They produce music during limited spare hours alongside their day jobs. Spending entire evenings dragging playheads across long recordings quickly drains creative energy. \textbf{Sample Browser} resolves this issue by indexing, tagging, and recommending music at the passage level (4 to 30 seconds), jumping creators straight to the cue points they need.

The core objective of this research is to design and validate the conceptual architecture of \textbf{Sample Browser}. The platform interfaces with public embedded audio via the YouTube IFrame API to deliver timestamped passage recommendations, avoiding media hosting costs and copyright risks. Guided by the Problem-Based Learning (PBL) framework, we grounded our engineering decisions in Music Information Retrieval (MIR) literature, creator surveys, and platform policy audits---specifically Spotify's 2024--2026 developer restrictions on audio feature endpoints~\cite{ref:spotifydev2024, ref:spotifydev2026}.

The following chapters map out our research trajectory and conceptual design. Chapter 2 details the problem analysis, documents our team's early methodological errors and iterative pivot, presents our solution architecture, establishes target personas with testable validation hypotheses (H1--H4), and benchmarks market alternatives alongside a strategic SWOT matrix. Chapter 3 concludes the report with our roadmap for Midterm 2.
